\documentclass[12pt]{article}
\usepackage[a4paper,margin=1in]{geometry}
\usepackage[T1]{fontenc}
\usepackage{lmodern,microtype}
\usepackage{amsmath,amssymb,amsthm,mathtools}
\usepackage{booktabs,array,enumitem,xcolor}
\usepackage[numbers,sort&compress]{natbib}
\usepackage[colorlinks=true,linkcolor=blue!55!black,citecolor=blue!55!black,urlcolor=blue!55!black]{hyperref}
\linespread{1.04}

\newtheorem{theorem}{Theorem}[section]
\newtheorem{proposition}[theorem]{Proposition}
\newtheorem{corollary}[theorem]{Corollary}
\theoremstyle{remark}
\newtheorem{remark}[theorem]{Remark}
\newcommand{\cH}{\mathcal H}
\newcommand{\cK}{\mathcal K}
\newcommand{\Ran}{\operatorname{Ran}}
\newcommand{\rank}{\operatorname{rank}}
\newcommand{\norm}[1]{\left\lVert#1\right\rVert}

\title{A Cyclic-Cloud Riesz-Shell Theorem\\
\large Explicit Root Geometry, Directed Schur Budgets, and Rank-Two Persistence}
\author{Liang Wang}
\date{July 2026}

\begin{document}
\maketitle

\begin{abstract}
The double-cycle model gives an exact finite cloud determinant, but a physical
application requires its roots to persist as Riesz shells of a generally
nonnormal transfer operator.  We specialize the RH-162--166 packet-to-Riesz
architecture to the normal cyclic seed and obtain completely explicit contour
geometry.

Let $L\ge3$, let $\rho>0$, and let
$K_L$ be the doubled reduced cycle with roots
$\zeta_k=\rho e^{2\pi ik/L}$, $1\le k<L$, each of multiplicity two.  The
half-spacing is
\[
 s_L=\rho\sin(\pi/L).
\]
Embed $K_L$ isometrically into a Hilbert space and write a candidate physical
operator in packet/complement blocks
\[
 A=\begin{pmatrix}K_L+E&B\\C&D\end{pmatrix}.
\]
Choose circles $\Gamma_k=\partial D(\zeta_k,\delta)$ with
$\delta<s_L$.  If $\norm E=\varepsilon<\delta$, the complement disks contain
no spectrum, their contour resolvent is bounded by $d$, and
\[
 \frac{d\norm B\norm C}{\delta-\varepsilon}<1,
\]
then every $\Gamma_k$ remains in the resolvent of $A$ and encloses a Riesz
projection of rank exactly two.  The proof uses an explicit packet resolvent
bound $(\delta-\varepsilon)^{-1}$ and a coupling homotopy controlled by the
directed Schur product.

A 192-matrix complex audit covers 1,248 root contours for
$L=4,5,6,8,10,12$.  Every admissible contour encloses rank two, with zero
certificate or rank failures.  The theorem is conditional: no physical
embedding, perturbation ball, complement resolvent, or all-level uniform
margin is supplied.  It converts a future physical cycle fit into explicit
Riesz obligations but does not close interface R or Gate A.
\end{abstract}

\section{Why the cyclic model helps the Riesz problem}

The general physical interface R of RH-171 requires an ambient realization,
validated block data, uniform Schur margins, and shellwise transport
\citep{WangRH171}.  For an arbitrary packet block, even choosing contours is
difficult because a nonnormal spectral gap does not control the resolvent.

The cyclic seed removes that ambiguity.  It is finite and normal, its roots
are explicit, and the nearest-neighbor separation is known exactly.  A
physical theorem would still need to show that some transfer-space block is
close to the seed, but once such closeness is available the packet resolvent
budget is closed form.

\section{The doubled reduced cycle}

Fix $L\ge3$ and $\rho>0$.  Let
\begin{equation}\label{eq:roots}
 \zeta_k=\rho e^{2\pi ik/L},
 \qquad 1\le k\le L-1.
\end{equation}
Let $K_L$ be the normal diagonalizable operator whose eigenvalue multiset
contains two copies of every $\zeta_k$.  Equivalently, $K_L$ is the direct
sum of two reduced length-$L$ cycles scaled by $\rho$.

The chord distance between adjacent roots is
\[
 |\zeta_{k+1}-\zeta_k|=2\rho\sin(\pi/L).
\]
Define the half-spacing
\begin{equation}\label{eq:spacing}
 s_L=\rho\sin(\pi/L).
\end{equation}
For every $0<\delta<s_L$, the closed disks
$\overline D(\zeta_k,\delta)$ are pairwise disjoint.

\begin{proposition}[Normal packet resolvent]\label{prop:normal-resolvent}
On $\Gamma_k=\partial D(\zeta_k,\delta)$,
\begin{equation}\label{eq:seed-resolvent}
 \norm{(z-K_L)^{-1}}=\frac1\delta.
\end{equation}
If $\norm E\le\varepsilon<\delta$, then
\begin{equation}\label{eq:perturbed-packet-resolvent}
 \norm{(z-K_L-E)^{-1}}
 \le\frac1{\delta-\varepsilon}
 =:a.
\end{equation}
\end{proposition}

\begin{proof}
For a normal finite matrix the resolvent norm is the reciprocal of the
distance to its spectrum.  The contour has distance $\delta$ from its central
root and greater distance from every other root.  The perturbation estimate
follows from
\[
 z-K_L-E=(I-E(z-K_L)^{-1})(z-K_L)
\]
and the Neumann series.
\end{proof}

\section{Physical block form}

Let $J:\mathbb C^{2(L-1)}\to\cH$ be an isometry and let $P=JJ^*$,
$Q=I-P$.  In the corresponding orthogonal decomposition, write
\begin{equation}\label{eq:block}
 A=\begin{pmatrix}
 K_L+E&B\\
 C&D
 \end{pmatrix}.
\end{equation}
Here $E$ is the packet-block modeling error, while $B$ and $C$ are the two
directed couplings.  They need not be adjoints.

Assume for every $k$ that
\begin{equation}\label{eq:complement-free}
 \overline D(\zeta_k,\delta)\subset\rho(D),
 \qquad
 \sup_{z\in\Gamma_k}\norm{(z-D)^{-1}}\le d.
\end{equation}
The disk condition ensures that the baseline complement contributes no
enclosed spectral rank.  A boundary inverse estimate alone would be
insufficient: complement eigenvalues could lie inside the contour.

\section{The cyclic Riesz-shell theorem}

\begin{theorem}[Rank-two cyclic shell persistence]\label{thm:riesz}
Assume
\begin{equation}\label{eq:geometry-assumptions}
 0<\delta<s_L,
 \qquad
 \norm E\le\varepsilon<\delta,
\end{equation}
and \eqref{eq:complement-free}.  If
\begin{equation}\label{eq:schur}
 \kappa:=\frac{d\norm B\norm C}{\delta-\varepsilon}<1,
\end{equation}
then every $\Gamma_k$ lies in $\rho(A)$ and its Riesz projection
\begin{equation}\label{eq:riesz-projection}
 \Pi_k=\frac1{2\pi i}\int_{\Gamma_k}(z-A)^{-1}\,dz
\end{equation}
has rank two.
\end{theorem}

\begin{proof}
Consider the homotopy
\[
 A_s=\begin{pmatrix}
 K_L+sE&sB\\
 sC&D
 \end{pmatrix},
 \qquad 0\le s\le1.
\]
For $z\in\Gamma_k$, Proposition~\ref{prop:normal-resolvent} gives
\[
 \norm{(z-K_L-sE)^{-1}}
 \le\frac1{\delta-s\varepsilon}
 \le\frac1{\delta-\varepsilon}=a.
\]
The upper-left Schur complement after eliminating $z-D$ is
\[
 z-K_L-sE-s^2B(z-D)^{-1}C.
\]
Relative to $z-K_L-sE$, the feedback norm is at most
\[
 a\,s^2\norm B\,d\,\norm C\le\kappa<1.
\]
Hence the Schur complement and therefore $z-A_s$ are invertible for every
$s$ and every point of the contour.

Riesz projections vary continuously in operator norm along a resolvent-
preserving homotopy, so their finite rank is constant \citep{Kato1995}.  At
$s=0$, the operator is $K_L\oplus D$.  The disk contains exactly one doubled
cycle root and, by \eqref{eq:complement-free}, no complement spectrum.
Therefore the enclosed rank is two at $s=0$ and at $s=1$.
\end{proof}

\section{Full resolvent and projector-distance budgets}

The Schur argument also gives an explicit full resolvent envelope.  Put
$a=(\delta-\varepsilon)^{-1}$,
$b=\norm B$, and $c=\norm C$.  On a root contour, let
\[
 \kappa=adbc<1.
\]
The inverse of the upper-left Schur complement has norm at most
$a/(1-\kappa)$.  The standard block inverse formula therefore yields
\begin{align}
 \norm{R_{11}(z)}&\le\frac{a}{1-\kappa},
 \label{eq:r11}\\
 \norm{R_{12}(z)}&\le\frac{abd}{1-\kappa},
 \label{eq:r12}\\
 \norm{R_{21}(z)}&\le\frac{dca}{1-\kappa},
 \label{eq:r21}\\
 \norm{R_{22}(z)}&\le\frac{d}{1-\kappa}.
 \label{eq:r22}
\end{align}
For example, the sum
\begin{equation}\label{eq:full-resolvent}
 M=\frac{a+abd+dca+d}{1-\kappa}
\end{equation}
is a conservative bound for the full block operator resolvent under the
direct-sum one-norm.

\begin{proposition}[Quantitative Riesz displacement]
\label{prop:projector-distance}
Along the homotopy in Theorem~\ref{thm:riesz}, suppose the full contour
resolvent is uniformly bounded by $M$.  Then
\begin{equation}\label{eq:projector-distance}
 \norm{\Pi_k(A)-\Pi_k(K_L\oplus D)}
 \le \delta M^2(\varepsilon+b+c).
\end{equation}
If the right side is less than one, the physical rank-two shell is a graph
over the corresponding doubled cycle eigenspace.
\end{proposition}

\begin{proof}
Differentiate the Riesz projection along $A_s$:
\[
 \frac{d\Pi_k(s)}{ds}
 =\frac1{2\pi i}\int_{\Gamma_k}
 (z-A_s)^{-1}A_s'(z-A_s)^{-1}\,dz.
\]
The derivative block has norm at most $\varepsilon+b+c$.  Since the contour
length is $2\pi\delta$, integration over $s\in[0,1]$ gives
\eqref{eq:projector-distance}.  Projection distance below one gives the usual
invertible graph map \citep{Kato1995}.
\end{proof}

These bounds turn RH-180 into a direct validation target: a physical audit
can report not only rank preservation but also graph distance and directional
off-diagonal resolvent blocks.

\section{Asymptotic scaling forced by growing cycle length}

For large $L$,
\begin{equation}\label{eq:spacing-asymptotic}
 s_L=\rho\sin(\pi/L)\sim\frac{\pi\rho}{L}.
\end{equation}
Choose a fixed fraction $\delta=\theta s_L$ and require
$\varepsilon\le\gamma\delta$ with $0<\gamma<1$.  Then
\[
 a=\frac1{\delta-\varepsilon}
 \in\left[\frac1{\theta s_L},
 \frac1{(1-\gamma)\theta s_L}\right]
 \asymp L.
\]
If the complement resolvent bound $d$ remains $O(1)$, the Schur condition
forces
\begin{equation}\label{eq:coupling-scaling}
 \norm B\norm C=O(L^{-1}).
\end{equation}
Thus bounded nonvanishing two-way coupling cannot support individually
resolved shells of unbounded cycle length.  At least one of the following is
needed: improving packet approximation, decaying directed feedback,
complement resolvent improvement, or grouping several nearby roots into a
coarser shell.

This scaling law is a genuine constraint on any all-level cyclic route.  The
finite RH-180 audit uses fixed small lengths and therefore tests the theorem,
not the required $L\to\infty$ decay.

\begin{remark}[Directed product, not a symmetric norm sum]
The criterion uses
$a\,d\,\norm B\,\norm C<1$.  A large coupling in one direction may be
harmless if the return coupling is small.  Replacing the product by a
symmetric requirement on $\norm B+\norm C$ discards this nonnormal structure,
as emphasized in RH-163 \citep{WangRH163}.
\end{remark}

\section{Consequences for the cloud determinant}

The Riesz projections $\Pi_k$ are pairwise annihilating because their
contours are disjoint.  Their sum
\[
 \Pi_{\rm cyc}=\sum_{k=1}^{L-1}\Pi_k
\]
has rank $2(L-1)$.  On its invariant range, the physical finite cloud factor
is
\begin{equation}\label{eq:physical-factor}
 C_A(w)=\det_{\Ran\Pi_{\rm cyc}}
 (I-wA|_{\Ran\Pi_{\rm cyc}}).
\end{equation}
The theorem guarantees its degree and shell multiplicities.  It does not make
$C_A$ exactly equal to the geometric polynomial when $E,B,C$ are nonzero.
That stronger identification requires root-location or coefficient bounds.

\begin{corollary}[Exact model limit]
If $E=B=C=0$, then $\Pi_{\rm cyc}=P$ and
\[
 C_A(w)=\Pi_{L-1}(\rho w)^2.
\]
For small nonzero data satisfying Theorem~\ref{thm:riesz}, the divisor splits
into the same $L-1$ rank-two shells but may move within the contours.
\end{corollary}

This distinction separates physical interface R (existence and rank of Riesz
shells) from physical interface Q (identification of their determinant
coefficients).

\section{Omission mechanisms}

Each hypothesis has a distinct role.
\begin{description}[leftmargin=3.5em,style=nextline]
\item[Root spacing.] If $\delta\ge s_L$, adjacent doubled roots can share a
contour, so individual rank-two shells are not defined.
\item[Packet perturbation.] If $\varepsilon\ge\delta$, a packet eigenvalue
can cross its contour even with zero complement coupling.
\item[Complement exclusion.] A bounded contour inverse does not rule out
complement eigenvalues inside the disk; rank two can fail at the baseline.
\item[Directed Schur margin.] In scalar blocks the contour determinant can
vanish exactly when the feedback product reaches one.
\item[Ambient isometry.] Without a target-independent $J$, the decomposition
\eqref{eq:block} is not a physical statement and can be fitted after seeing
the desired spectrum.
\end{description}
Thus the theorem is explicit but not self-instantiating.

\section{Finite matrix audit}

The audit uses
\[
 L=4,5,6,8,10,12,
 \qquad \rho=\lambda^{-1},
 \qquad \delta=0.35s_L,
 \qquad \varepsilon=0.15\delta.
\]
For each length, 32 independent complex perturbations and directed couplings
are generated.  The complement has three normal eigenvalues on radius
$2.5\rho$, and the coupling norms are chosen so that the certified Schur
product equals $0.20$.

This gives 192 full matrices and 1,248 individual root contours.  Every
budget is admissible, every contour encloses exactly two numerical
eigenvalues, and the audit records zero rank failures and zero certificate
failures.  These tests validate the formula and counting implementation; they
are not samples of the physical transfer operator.

\section{Physical obligations and next frontier}

Theorem~\ref{thm:riesz} provides a finite conditional route from a cyclic
model to Riesz shells.  To apply it physically one must still:
\begin{enumerate}[leftmargin=2em]
\item construct an isometry $J$ from a data-derived cycle into the noisy
transfer/determinant space;
\item validate $\norm E$, both coupling norms, and complement resolvents with
outward bounds;
\item prove an eventual uniform margin as $L$ grows and root spacing shrinks
like $\pi\rho/L$;
\item transport every fixed shell in common coordinates across scales;
\item separately identify the resulting cloud factor and directed marks.
\end{enumerate}
The shrinking spacing is important: a fixed absolute modeling error cannot
support arbitrarily long cycles.  The physical approximation must improve at
least on the $L^{-1}$ contour scale.

No physical input above is proved here.  Interface R, the cloud ledger Q, the
Schatten complement U, normalization Z, directed limit T, macro Gate A, and
all Hilbert--Polya or Riemann-hypothesis conclusions remain open.

\bibliographystyle{plainnat}
\bibliography{references}
\end{document}
